\textbf{Amaç} Elma bahçelerinde çalışan toplama robotu filolarının binlerce ağaç üzerinde
rotalanması ve bir robot her döndüğünde yeniden planlanması gerekir, çünkü ağaç başına olgun meyve
miktarı ancak toplama sırasında belirlenir. Mevcut min--maks rotalama çözücüleri
dakikalardan saatlere uzanan bütçelerle değerlendirilmektedir. Bu çalışmada, bahçe ölçeğinde iyi filo
planlarının ne kadar hızlı elde edilebileceği ve gerçekçi bir hasat senaryosunda hızlı yeniden planlamanın
ne kadar yarar sağladığı araştırılmıştır.

\textbf{Yöntem} Önce rota, sonra küme yaklaşımına dayalı bir planlayıcı olan PSB, Greedy-Edge dev turunu
min--maks amacına göre optimal biçimde bölmekte, rotaları dengelemekte ve her an sonuç veren (anytime)
bir aramayla iyileştirmektedir; yöntem ayrıca kasa kapasiteli seferleri ve ağaç konumlarından kestirilen
bahçe eksenini kapsayacak biçimde genişletilmiştir. PSB, tek çekirdek üzerinde LKH-3 ve literatürde öne çıkan
dört kodla, gerçek bahçe sınırları da dâhil \InstAll{} örnek üzerinde ve parametreleri
literatürden alınan, kablolu drone hasat platformlarına ilişkin bir benzetimde karşılaştırılmıştır.

\textbf{Bulgular} PSB ilk çözümü bir saniye içinde döndürmüş, 10\,s sonunda MILS'i örneklerin
\SotaMilsWinTenAll\% kadarında geride bırakmış ve 5000'den fazla hedef içeren örneklerde
600\,s sonunda onunla başa baş kalmıştır; buna karşın orta büyüklükteki örneklerde, birkaç dakikalık
süre bulunduğunda MILS daha iyi sonuç vermiştir. Hasat benzetiminde PSB-MT/R ile yeniden planlama, ilk toplama
geçişini statik sıra bantlarına kıyasla koşuların \HgFirstRowblockWin\% kadarında kısaltmış,
bahçe ekseninin göz ardı edilmesi ise görevleri
\HgFramePsbAxisZeroAbs\% uzatmıştır.

\textbf{Sonuç} Hızlı ve bahçe eksenine duyarlı yeniden planlama, benzetilen bahçe ölçekli
toplama filolarında görev süresini belirleyen önemli etkenlerdendir.
